\documentclass[letterpaper, onecolumn, draftclsnofoot, 10pt, compsoc]{IEEEtran}
\usepackage{graphicx}
\usepackage[hyphens]{url}
\usepackage{setspace}

\usepackage{geometry}
\geometry{textheight=9.5in, textwidth=7in}

% 1. Fill in these details
\def \CapstoneTeamName{		The Polycules}
\def \CapstoneTeamNumber{		63}
\def \GroupMemberOne{			Corynna Park}
\def \GroupMemberTwo{			Joshua Lioy}
\def \GroupMemberThree{			Jackson Wells}
\def \CapstoneProjectName{		Molecules in 3D?! And in color!? That I can hold in my hand? No way!!!}
\def \CapstoneSponsorCompany{	Oregon State University College of Science Department of Biochemistry and Biophysics}
\def \CapstoneSponsorPerson{		Dr. Victor Hsu}

% 2. Uncomment the appropriate line below so that the document type works
\def \DocType{	%Problem Statement
				Requirements Document
				%Technology Review
				%Design Document
				%Progress Report
				}
			
\newcommand{\NameSigPair}[1]{\par
\makebox[2.75in][r]{#1} \hfil 	\makebox[3.25in]{\makebox[2.25in]{\hrulefill} \hfill		\makebox[.75in]{\hrulefill}}
\par\vspace{-12pt} \textit{\tiny\noindent
\makebox[2.75in]{} \hfil		\makebox[3.25in]{\makebox[2.25in][r]{Signature} \hfill	\makebox[.75in][r]{Date}}}}
% 3. If the document is not to be signed, uncomment the RENEWcommand below
%\renewcommand{\NameSigPair}[1]{#1}

%%%%%%%%%%%%%%%%%%%%%%%%%%%%%%%%%%%%%%%
\begin{document}
\begin{titlepage}
    \pagenumbering{gobble}
    \begin{singlespace}
    	\includegraphics[height=4cm]{coe_v_spot1}
        \hfill 
        % 4. If you have a logo, use this includegraphics command to put it on the coversheet.
        %\includegraphics[height=4cm]{CompanyLogo}   
        \par\vspace{.2in}
        \centering
        \scshape{
            \huge CS Capstone \DocType \par
            {\large\today}\par
            \vspace{.5in}
            \textbf{\Huge\CapstoneProjectName}\par
            \vfill
            {\large Prepared for}\par
            \Huge \CapstoneSponsorCompany\par
            \vspace{5pt}
            {\Large\NameSigPair{\CapstoneSponsorPerson}\par}
            {\large Prepared by }\par
            Group\CapstoneTeamNumber\par
            % 5. comment out the line below this one if you do not wish to name your team
            \CapstoneTeamName\par 
            \vspace{5pt}
            {\Large
                \NameSigPair{\GroupMemberOne}\par
                \NameSigPair{\GroupMemberTwo}\par
                \NameSigPair{\GroupMemberThree}\par
            }
            \vspace{20pt}
        }
        \begin{abstract}
        % 6. Fill in your abstract    
        \end{abstract}     
    \end{singlespace}
\end{titlepage}
\newpage
\pagenumbering{arabic}
\tableofcontents
% 7. uncomment this (if applicable). Consider adding a page break.
%\listoffigures
%\listoftables
\clearpage

% 8. now you write!
\section{Introduction}

\subsection{Purpose} %jackson
This document is intended to give a detailed overview the requirements of our polychromatic 3D printing project.
Contained in this document will be the information required for an individual to gauge the completeness of our project.
Individuals looking to reproduce the project will be able to use this document, to verify that they created an identical product.
Audiences concerned with this document include owners of multi-filament interfaces and single extruder printers, anyone with access to these tools or anyone with an interest in creating polychromatic 3D objects.

\subsection{Scope} %Josh
Upon completion of our polychromatic 3d printing project we will have produce a polychromatic 3d printing work-flow that is comprised of various open source and proprietary components. This work-flow will allow a user who has basic knowledge of 3d printing to print a polychromatic object just as easily as a monochromatic object. The work-flow will be able to take monochromatic 3d object files, and convert them into polychromatic 3d object files that can be printed with the appropriate hardware.\par
Initially we hope that the work-flow will allow for faculty and staff in the Oregon State University College of Science Department of Biochemistry and Biophysics to efficiently create polychromatic scale representations of molecules for educational and instructional purposes. We see this work-flow as being a way to make better use of the technology that the department already has available as well as a potential teaching tool in the future.\par
Beyond the specific application of the work-flow to the Department of Biochemistry and Biophysics we hope that we will be able to adapt our work-flow to any polychromatic 3d printing application.\par

\subsection{Definitions, Acronyms and Abbreviations} %Josh
  \begin{itemize}
  	\item Slicing Software: A program that takes a 3d object file and translates it into individual layers that can be interpreted by a 3d printer.
    \item G-code: A numerical control programming language that is used in computer-aided manufacturing and 3d printing.
    \item STL: Short for Stereo-Lithography, STL is a file format that is widely used in 3d printing to convert object data into G-code.
  \end{itemize}

\subsection{References} %Josh
%We will need to decide if we want the references in an appendix or here.

\subsection{Overview} %Corynna
The remainder of this document will discuss the general factors that affect the product and its requirements as well as give detailed descriptions of all product requirements. 
Consequently, it will be divided into two main sections, and each section will be further split up into subsections.
Section 2 will give the background information for the product such that the product requirements will be easier to understand.
Section 3 will describe the product requirements such that the completeness of the project will be clearly defined by the fulfillment of these requirements.
All requirements will be feasible and testable.

\section{Overall Description}
Summary of section 3

\subsection{Product Perspective} %Jackson
Currently, no method exists for producing objects with multiple colors.
The polychromatic 3D printing pipeline will be the first of its kind, consisting of multiple open-source software packages. 
The pipeline will work together with a filament interface and 3D printer, to print 3D objects. 
The product will be interfaced with via desktop computer, without put being sent to a USB drive.

\subsection{Product Functions} %Jackson
\begin{itemize}
	\item Produce multi-colored structurally viable 3D objects from biological data.
    \item Produce multi-colored structurally viable 3D objects.
    \item Process, annotate, and convert provided object files for use by printer and multi-filament interface.
\end{itemize}

\subsection{User Characteristics} %Corynna
At a minimum, users should have an interest in polychromatic 3D printing along with basic knowledge of how 3D printers work and how to operate them. 
While our project originated from the Biochemistry department, users will not need to have knowlege of this field, as the product will be intended for a wider audience.
Users should be of the technical expertise to where they are able to read, comprehend, and follow a technical workflow which includes instruction to use certain software and generate specific files. 
Users should own or have access to multi-filament interfaces and single extruder printers. 

\subsection{Constraints} %Jackson
\begin{itemize}
	\item Software used is open-source and available free to the public.
    \item Software runs on Mac and Windows OS.
    \item Software used must be able to output object files in GCODE for use by the printer and interface.
\end{itemize}

\subsection{Assumption and Dependencies} %Corynna
The product requirements will be based on the following assumptions and dependencies: 
\begin{itemize}
  \item There is preexisting software that allows the generation of files to print monochromatic objects in 3D.
  \item There exists free open source software that allows the generation of multiple files to represent a polychromatic object in 3D.
  \item The file(s) representing the 3D object will be fed to the 3D printer through USB. 
  \item The user has access to a multi-filament interface and single extruder printer.
\end{itemize}

\section{Specific Requirements} %everyone

\cite{8colorprint}

\bibliography{references}{}
\bibliographystyle{ieeetr}

\end{document}
